\chapter{Tái lập và chỉ mục nguồn}
\label{app:tai-lap}

Phụ lục này mô tả cách dựng lại môi trường, sinh lại bảng/hình, và chạy lại từng thực nghiệm.
Mục tiêu là để một người khác có thể kiểm tra một con số trong báo cáo mà không cần hỏi nhóm.

\section{Môi trường}
\paragraph{Công cụ.} Dự án dùng \texttt{uv} để quản lý Python và thư viện. Cài đặt:
\begin{lstlisting}[numbers=none]
git clone https://github.com/SGU-ML25/CaroGame
cd CaroGame
uv sync --group dev --extra nn --extra ui
uv run pytest
\end{lstlisting}
Lệnh cuối phải in ``258 passed, 1 skipped'' tại thời điểm viết. Nếu số test khác, có thể mã
nguồn đã thay đổi so với bản dùng trong báo cáo; kiểm tra \texttt{meta.json} của từng run để
biết commit tương ứng.

\paragraph{Lõi Cython.} Alpha-Beta và PUCT có bản Cython tùy chọn:
\begin{lstlisting}[numbers=none]
uv pip install cython setuptools
uv run python setup_native.py build_ext --inplace
\end{lstlisting}
Nếu chưa build, code tự quay về Python. Kết quả lý thuyết giống nhau; chỉ thời gian khác. Các
bảng thời gian trong báo cáo ghi rõ backend.

\paragraph{Model.} Trọng số L5/L6 nằm trên Hugging Face, ghim trong \texttt{configs/models.yaml}.
L8 dùng checkpoint cục bộ \repoFile{runs/train/alphazero/20261124-160615\_gpu-native/checkpoints/best.pt}.
Nếu checkpoint không có, có thể tải L5/L6 bằng:
\begin{lstlisting}[numbers=none]
uv run caro-hub pull l5 l6
\end{lstlisting}
L8 chưa được push lên Hugging Face tại thời điểm viết; muốn tái lập đúng kết quả L8 cần file
\texttt{best.pt} của run gốc, hoặc chạy lại 50 vòng huấn luyện (khoảng 12 giờ trên GPU laptop).

\section{Bản đồ thư mục}
\begin{center}\small
\begin{tabularx}{\textwidth}{@{}lX@{}}
\toprule
\textbf{Đường dẫn} & \textbf{Nội dung} \\
\midrule
\repoFile{caro/engine/} & Bitboard, CaroGame, Zobrist, đối xứng D4 \\
\repoFile{caro/agents/} & Tám cấp tác tử, giao diện chung, PUCT \\
\repoFile{caro/nn/} & Đặc trưng, mạng ResNet policy--value, dataset \\
\repoFile{caro/training/} & Sinh dữ liệu, SL, REINFORCE, vòng lặp AlphaZero \\
\repoFile{caro/arena/} & Trận, giải, khai cuộc, thống kê, benchmark \\
\repoFile{caro/ui/} & FastAPI và giao diện web \\
\repoFile{configs/} & Toàn bộ tham số thực nghiệm \\
\repoFile{runs/} & Log thô của mọi lần chạy \\
\repoFile{report/scripts/} & Script sinh bảng, hình, phụ lục \\
\repoFile{report/ban-mo-rong/} & Nội dung báo cáo mở rộng và bảng sinh tự động \\
\bottomrule
\end{tabularx}
\end{center}

\section{Sinh lại bảng và hình}
Từ gốc repo, chạy:
\begin{lstlisting}[numbers=none]
uv run python report/scripts/phan-tich-tong-hop.py
uv run python report/scripts/phan-tich-hoc.py
uv run python report/scripts/phu-luc-chon-loc.py
uv run python report/scripts/ho-so-van-moi.py
cd report && latexmk -xelatex bao-cao-caro.tex
\end{lstlisting}
Các script chỉ đọc log trong \repoFile{runs/} và ghi vào \repoFile{report/ban-mo-rong/}. Chúng
không sửa log gốc. Nếu một bảng không sinh được, kiểm tra đường dẫn run tương ứng còn tồn tại
không; báo cáo dùng đường dẫn tuyệt đối tương đối so với gốc repo.

\section{Chạy lại từng thực nghiệm}
\paragraph{Giải tám tác tử (RR8).}
\begin{lstlisting}[numbers=none]
uv run python -m caro.arena.tournament experiments/rr8_mode_b
\end{lstlisting}
Cấu hình dùng 50 khai cuộc, 4 worker mặc định. Trong báo cáo, run được chạy với
\texttt{arena.workers=1} để không ảnh hưởng máy đang dùng; muốn nhanh hơn có thể tăng worker.

\paragraph{Tự đấu từng cấp (E3).}
\begin{lstlisting}[numbers=none]
uv run python experiments/e3_selfplay.py --workers 4
\end{lstlisting}
Mỗi cấp chơi 50 khai cuộc × 2 màu. Có thể giới hạn cấp bằng \texttt{--levels}.

\paragraph{Các ablation.}
\begin{lstlisting}[numbers=none]
uv run python -m caro.arena.bench experiments/a1_ab_pruning
uv run python -m caro.arena.bench experiments/a2_ordering_tt
uv run python -m caro.arena.bench experiments/a7_eval_weights
uv run python -m caro.arena.tournament experiments/a7b_head_to_head
uv run python -m caro.arena.tournament experiments/a8b_root_tactics
\end{lstlisting}

\paragraph{Huấn luyện.}
\begin{lstlisting}[numbers=none]
uv run python -m caro.training.gen_expert_data training/gen_expert_data_v2
uv run python -m caro.training.train_sl training/sl_v3
uv run python -m caro.training.train_pg training/pg
uv run python -m caro.training.alphazero_loop training/alphazero
\end{lstlisting}

\section{Chỉ mục nguồn số liệu}
\begin{small}
\begin{longtable}{@{}p{.24\textwidth} p{.68\textwidth}@{}}
\toprule
\textbf{Phép đo} & \textbf{Đường dẫn log} \\
\midrule
RR8 & \repoFile{runs/eval/tournament/rr8\_mode\_b/latest/} \\
E3 & \repoFile{runs/eval/tournament/e3\_selfplay/20261125-232320/} và \repoFile{20261125-233229/} \\
A1 & \repoFile{runs/bench/a1\_ab\_pruning/20261029-200321/} \\
A2 & \repoFile{runs/bench/a2\_ordering\_tt/20261029-212610/} \\
A7 & \repoFile{runs/bench/a7\_eval\_weights/20261029-223250/} \\
A3 & \repoFile{runs/tournament/a3\_mcts\_sims/20261114-232353/} \\
A4 & \repoFile{runs/tournament/a4\_rollout/20261114-234556/} \\
A5 & \repoFile{runs/eval/a5\_puct/results.csv} \\
Arena cũ & \repoFile{runs/arena\_matrix/games.csv} và \repoFile{ranking.csv} \\
SL v1/v2/v3 & \repoFile{runs/train/sl\_policy*/metrics.json} \\
RL v2/v3 & \repoFile{runs/train/pg\_policy/20261117-224555\_v2\_init/} và \repoFile{20261118-125103\_v3\_init/} \\
L8 & \repoFile{runs/train/alphazero/20261124-160615\_gpu-native/} \\
L8--L3 & \repoFile{runs/eval/tournament/l8\_vs\_l3/20261125-134954\_round37/} \\
\bottomrule
\end{longtable}
\end{small}

\section{Đối chiếu checkpoint}
\begin{center}\small
\begin{tabularx}{\textwidth}{@{}l l X@{}}
\toprule
\textbf{Tên} & \textbf{Loại} & \textbf{Ghi chú} \\
\midrule
\texttt{model:l5} & SL v3 & Trọng số dùng cho L5 và L7; revision ghim trong \texttt{models.yaml} \\
\texttt{model:l6} & RL v3 & Trọng số dùng cho L6; revision ghim trong \texttt{models.yaml} \\
\texttt{model:l8} & L8 & Revision null tại thời điểm viết; dùng checkpoint cục bộ \\
\texttt{sl\_policy/latest.pt} & SL v1 & Bản cũ, nhãn value tuyệt đối; chỉ dùng để so lịch sử \\
\texttt{sl\_policy\_v2/latest.pt} & SL v2 & Bản cũ, nhãn value tuyệt đối; chỉ dùng để so lịch sử \\
\texttt{sl\_policy\_v3/latest.pt} & SL v3 & Bản đã sửa quy ước; dùng làm init cho RL và PUCT \\
\bottomrule
\end{tabularx}
\end{center}

\section{Kiểm tra một con số}
Muốn kiểm tra một dòng trong bảng, làm theo ba bước. Thứ nhất, mở \texttt{config.yaml} của run
tương ứng để xem ngân sách, seed và danh sách tác tử. Thứ hai, mở \texttt{games.jsonl} và lọc
theo \texttt{agent\_x}, \texttt{agent\_o}; mỗi dòng là một ván, \texttt{result} theo góc nhìn
X. Thứ ba, dùng script phân tích hoặc tự tính $(W+0.5D)/n$ để đối chiếu. Nếu số khác, kiểm tra
xem bảng có dùng bootstrap theo cặp khai cuộc hay Wilson từng ván; hai phương pháp cho khoảng
khác nhau.

\section{Giới hạn của tái lập}
Không phải mọi thứ đều tái lập tuyệt đối. Thời gian chạy phụ thuộc máy và tải; kết quả PUCT có
thể phụ thuộc phiên bản PyTorch/CUDA; các run cũ có thể dùng code khác bản hiện tại. Báo cáo
ghi commit trong \texttt{meta.json} để người đọc biết. Các run dirty được đánh dấu và không
dùng làm bằng chứng chính khi có run sạch tương đương. Khi tái lập, nên so kết quả trong cùng
một môi trường, không so tuyệt đối với số trong báo cáo.